\section*{Chapitre \ref{ch:b3:surfaces-catalysis} --- Surfaces, adsorption et catalyse hétérogène}

\begin{solution}{exo:b3:surfaces-catalysis:1}
$K = \theta/((1 - \theta)p) = 0.40/(0.60 \times 2.0) = \qty{0.33}{kPa^{-1}}$~; sous \qty{10}{kPa}, $\theta = 3.33/4.33 = 0.77$.
\end{solution}

\begin{solution}{exo:b3:surfaces-catalysis:2}
$-20$~: \omterm{def:b3:surfaces-catalysis:physi-chemi}{physisorption} (l'ordre de grandeur d'une enthalpie de
condensation)~; $-150$~: \omterm{def:b3:surfaces-catalysis:physi-chemi}{chimisorption}, qui seule peut être
dissociative.
\end{solution}

\begin{solution}{exo:b3:surfaces-catalysis:3}
(100)~: un site sommet, deux sites pontés (quatre liaisons partagées
chacune entre deux atomes), un creux d'ordre quatre (quatre creux
partagés chacun entre quatre atomes). (111)~: un site sommet, trois
sites pontés, deux creux d'ordre trois (six triangles autour de chaque
atome, chacun partagé entre trois).
\end{solution}

\begin{solution}{exo:b3:surfaces-catalysis:4}
$\num{3.0e-6}/\num{1.5e-5} = \qty{0.20}{s^{-1}}$.
\end{solution}

\begin{solution}{exo:b3:surfaces-catalysis:5}
$\theta/(1 - \theta) = (Kp)^{1/2}$~: $0.111$ sous \qty{1.0}{kPa}, doublé sous \qty{4.0}{kPa}~: $0.222$, donc $\theta = 0.18$.
\end{solution}

\begin{solution}{exo:b3:surfaces-catalysis:6}
Dénominateur $1 + 10 + 2.5 = 13.5$~: $\theta_{\ce{CO}} = 0.74$,
$\theta_{\ce{O2}} = 0.19$, sites libres 0.07. La vitesse de
Langmuir--Hinshelwood, proportionnelle à $\theta_{\ce{CO}}\theta_{\ce{O2}}$,
est limitée par la rareté du dioxygène sur une surface encombrée de CO.
\end{solution}

\begin{solution}{exo:b3:surfaces-catalysis:7}
$\Delta_{\mathrm{ads}}H = -R\ln(8.0/1.0)/(1/300 - 1/350) = -8.314 \times 2.079/\num{4.76e-4} = \qty{-36}{kJ/mol}$.
\end{solution}

\begin{solution}{exo:b3:surfaces-catalysis:8}
$n_{\mathrm m} = 120/22\,414 = \qty{5.35e-3}{mol}$~; aire $= \num{5.35e-3} \times \num{6.022e23} \times \num{0.162e-18} =
\qty{522}{m^2.g^{-1}}$.
\end{solution}

\begin{solution}{exo:b3:surfaces-catalysis:9}
Langmuir--Hinshelwood~: les deux réactifs doivent être adsorbés sur des
sites voisins~; CO se fixe fortement et, à haute pression, ne laisse pas
de place au dioxygène.
\end{solution}

\begin{solution}{exo:b3:surfaces-catalysis:10}
Faible \omterm{def:b3:surfaces-catalysis:coverage}{taux de recouvrement}~: $100 - 60 = \qty{40}{kJ/mol}$~; fort \omterm{def:b3:surfaces-catalysis:coverage}{taux
de recouvrement}~: \qty{100}{kJ/mol}. Le tracé de $\ln r$
en fonction de $1/T$ est raide (pente $-100/R$) à basse température, où
la surface est pleine, et moins raide (pente $-40/R$) à haute
température, où elle se vide~: une courbe qui s'infléchit entre les
deux.
\end{solution}

\begin{solution}{exo:b3:surfaces-catalysis:11}
$4 \times 0.10 = 0.40$ des sites sont bloqués~: il reste 60\,\% de
l'activité pour un site, environ $0.60^2 = 36\,\%$ pour deux sites libres
adjacents.
\end{solution}

\begin{solution}{exo:b3:surfaces-catalysis:12}
Le métal intermédiaire (\qty{-250}{kJ/mol})~: le premier fixe l'azote
trop faiblement pour rompre facilement \ce{N2}, le troisième trop
fortement, si bien que les atomes d'azote restent sur la surface et la
bloquent.
\end{solution}

\begin{solution}{pb:b3:surfaces-catalysis:1}
\textbf{1.} Le diazote se condense vers \qty{77}{K} sous la pression
atmosphérique~: on atteint tout le domaine des pressions relatives, et
des multicouches se forment.
\textbf{2.} Au-dessous de 0.05, l'hétérogénéité de la surface, et
au-dessus de 0.30, la condensation dans les pores, violent les hypothèses
BET.
\textbf{3.} \num{1.278e-3}, \num{2.410e-3}, \num{3.453e-3}, \num{4.621e-3}, \num{5.659e-3}, \num{6.813e-3} \unit{g.cm^{-3}}.
\textbf{4.} $v_{\mathrm m} = 1/(0.02205 + 0.000180) = \qty{45.0}{cm^3.g^{-1}}$.
\textbf{5.} $c = 1 + 0.02205/0.000180 = 124$~: la première couche est
liée bien plus fortement que le liquide.
\textbf{6.} Elle est minuscule devant le produit de la pente par $x$, si
bien que son erreur relative est grande~; mais $v_{\mathrm m}$ dépend de
la somme de la pente et de l'ordonnée à l'origine, dominée par la pente.
\textbf{7.} $45.0/22\,414 = \qty{2.01e-3}{mol.g^{-1}}$.
\textbf{8.} $\num{2.01e-3} \times \num{6.022e23} \times \num{0.162e-18} = \qty{196}{m^2.g^{-1}}$.
\textbf{9.} Le support d'alumine~: le platine exposé (partie III)
représente environ \qty{0.9}{m^2.g^{-1}}.
\textbf{10.} Des pores de largeur moléculaire se remplissent entièrement
à très basse pression, et non couche par couche, si bien qu'aucune
\omterm{def:b3:surfaces-catalysis:coverage}{monocouche} ne peut être identifiée.
\textbf{11.} $2 \times 0.205/22\,414 = \qty{1.83e-5}{mol.g^{-1}}$.
\textbf{12.} $0.0100/195.08 = \qty{5.13e-5}{mol.g^{-1}}$.
\textbf{13.} $D = 1.83/5.13 = 0.357$.
\textbf{14.} $a^3/4 = (0.3923)^3/4 = \qty{0.01509}{nm^3}$.
\textbf{15.} $\frac{\sqrt3}{4}(0.3923)^2 = \qty{0.0666}{nm^2}$, la face
la plus dense~; (100) et (110) donnent 0.0770 et
\qty{0.1088}{nm^2}, et la moyenne des trois densités de sites donne
\qty{0.0807}{nm^2}.
\textbf{16.} Une sphère contient $\pi d^3/6v_{\mathrm{at}}$ atomes, dont $\pi d^2/a_{\mathrm{at}}$ en surface~:
$D = 6v_{\mathrm{at}}/(a_{\mathrm{at}}d)$.
\textbf{17.} $d = 6 \times 0.01509/(0.0807 \times 0.357) = \qty{3.1}{nm}$.
\textbf{18.} Un H par Pt de surface~; des particules sphériques d'une
seule taille (le résultat est une moyenne pondérée par la surface)~;
tout le platine réduit et accessible~; aucun débordement d'hydrogène
sur le support.
\textbf{19.} $\num{2.0e-5}/\num{1.83e-5} = \qty{1.1}{s^{-1}}$.
\textbf{20.} $\num{2.0e-5}/0.0100 = \qty{2.0e-3}{mol.g^{-1}.s^{-1}}$ par
gramme de platine.
\textbf{21.} $D = 1.12/10 = 0.112$ au lieu de 0.357~: un facteur 3.2.
\textbf{22.} Que la réaction est sensible à la structure~: ses \omterm{def:b3:complex-kinetics:enzyme}{sites
actifs} (arêtes, sommets, faces particulières) ne représentent pas une
fraction constante de la surface.
\textbf{23.} Elle compte les événements par \omterm{def:b3:complex-kinetics:enzyme}{site actif}, ce qui élimine
l'effet de la quantité de métal exposée.
\textbf{24.} \textbf{Diamètre moyen des particules de platine
$\approx \qty{3.1}{nm}$.}
\end{solution}
